

\section{What, Why, How}

\begin{frame}{Goal}
	\begin{itemize}
		\item Let's build a motion capture suite using microphones and speakers
		\item Think \href{https://www.lightningmaps.org/\#m=ses;t=3;s=0;o=0;b=;ts=0;z=5;y=35.9448;x=-87.045;d=2;dl=2;dc=0;}{real-time lightening map}
		\item Joint self-calibration and localization in real time for multi-emitter multi-receiver system
	\end{itemize}
\end{frame}

\begin{frame}{Environment}
	\begin{block}{Environment}
		\begin{itemize}
			\item Teensy 4.1 (600MHz Arm Cortex-M7) for perf tests and deployment
			\item PJRC Audio Board (two 2-channel boards = 4-channel audio)
			\item CMake (\href{https://github.com/newdigate/teensy-cmake-macros}{github teensy-cmake-macros})
			\item Arm64 Apple M2 for perf and unit tests
		\end{itemize}
	\end{block}
	
	\begin{block}{Constraints}
		\begin{itemize}
			\item Audio samples are \texttt{signed int16\_t}
			\item Audio channels are time aligned (phase locked)
			\item Audio is streaming into the Teensy at 44100 Hz
			\item Audio samples are grouped into "blocks" of 128 samples
			\item "Real time" means we have 2.9ms to process a block of 128 samples for each of the 4 input channels
			\item \textbf{No pre-existing FFT library for Arduino/Teensy that I liked}
		\end{itemize}
	\end{block}
\end{frame}


\begin{frame}{Approach}
	\begin{itemize}
		\item Use Time Difference of Arrival (TDOA) of a known chirp to estimate receiver/emitter geometry
		\item The more accurate the TDOA estimates, the more accurate  our receiver/emitter geometry estimates can be
		\item Typical CD Sampling Frequency (44100Hz) insufficient for naive algorithms ($ \leq 15$mm position error) - we can do better
		\item Use frequency domain methods to resolve TDOA beyond the microphone sampling rate
	\end{itemize}
\end{frame}


% \section{What is TDOA}

\begin{frame}{Time Difference of Arrival Geometry}
	\begin{figure}
	\centering
	\begin{tikzpicture}[scale=0.5,>=Latex]
		
		%----------------------------------------
		% Receiver locations
		%----------------------------------------
		\coordinate (R1) at (2.5,0);
		\coordinate (R2) at (-2.5,0);
		
		%----------------------------------------
		% Source
		%----------------------------------------
		\coordinate (S) at (3.6,3.0);
		
		%----------------------------------------
		% Compute distances
		%----------------------------------------
		\path let \p1 = ($(R1)-(S)$) in
		\pgfextra{\xdef\radius{veclen(\x1,\y1)}};
		
		\path let \p1 = ($(R2)-(S)$) in
		\pgfextra{\xdef\distSRtwo{veclen(\x1,\y1)}};
		
		%----------------------------------------
		% Hyperbola locus
		% Shifted to align with receiver geometry
		%----------------------------------------
		% Receiver locations
		\coordinate (R1) at (2.5,0);
		\coordinate (R2) at (-2.5,0);
		
		% Source location
		\coordinate (S) at (3.6,3.0);
		
		
		%----------------------------------------
		% Compute hyperbola coefficients
		%----------------------------------------
		
		% c = half receiver spacing
		\pgfmathsetmacro{\cval}{2.5}
		
		% Distances from source to receivers
		\pgfmathsetmacro{\distOne}{sqrt((3.6-2.5)^2+(3.0)^2)}
		\pgfmathsetmacro{\distTwo}{sqrt((3.6+2.5)^2+(3.0)^2)}
		
		\pgfmathsetmacro{\aval}{abs(\distTwo-\distOne)/2}
		
		% b = sqrt(c^2-a^2)
		\pgfmathsetmacro{\bval}{sqrt(\cval*\cval-\aval*\aval)}
		
		% midpoint of receivers
		\pgfmathsetmacro{\xcenter}{(2.5+(-2.5))/2}
		
		% Right branch of hyperbola
		\draw[gray,thin,samples=200,smooth,domain=-2:2]
		plot ({\xcenter+\aval*cosh(\x)}, {\bval*sinh(\x)});
		
		%\node[blue] at (1.0,2.2) {Hyperbola};
		
		%----------------------------------------
		% Circle centered at source through R1
		%----------------------------------------
		\draw[gray,thick] (S) circle (\radius);
		
		%----------------------------------------
		% Point where S-R2 line intersects circle
		%----------------------------------------
		\pgfmathsetmacro{\ratio}{\radius/\distSRtwo}
		\coordinate (P) at ($(S)!\ratio!(R2)$);
		
		%----------------------------------------
		% Radius d to nearer receiver
		%----------------------------------------
		\draw[very thick]
		(S)--(R1)
		node[midway,above right] {$d$};
		
		%----------------------------------------
		% Same distance toward R2
		%----------------------------------------
		\draw[very thick]
		(S)--(P);
		
		%----------------------------------------
		% Excess path length
		%----------------------------------------
		\draw[very thick,dashed,red]
		(P)--(R2)
		node[midway,below] {$\Delta d$};
		
		%----------------------------------------
		% Points
		%----------------------------------------
		\fill (R1) circle (2pt);
		\fill (R2) circle (2pt);
		\fill[red] (S) circle (2.5pt);
		
		\node[below] at (R1) {$R_1$};
		\node[below] at (R2) {$R_2$};
		\node[above left] at (S) {S};
		
	\end{tikzpicture}
	\caption{Example TDOA Locus. $S$ is a source, $R$ is a receiver pair, and $\Delta d$ is the difference in arrival times expressed as an equivalent distance. The half hyperbola shows the set of possible source locations for a given $\Delta d$. }
	\end{figure}
\end{frame}

\begin{frame}{Where We Are Going}
	\begin{figure}
		\centering
		\includegraphics[width=1.0 \textwidth]{../../python_trilateration_sandbox/tdoa/loci_from_truth_points_2d.png}
		\caption{Pairwise receiver-emitter TDOA locus (2D)}
	\end{figure}
\end{frame}


\begin{frame}{Approach - Mock Chirp}
	We will be looking for this chirp in microphone recordings
	\begin{figure}
		\centering
%		\includegraphics[width=0.6 \textwidth]{../../python_correlation_sandbox/modern/saved_figs/06_example_sinct.pdf}
		\includegraphics[width=0.6 \textwidth]{../../python_correlation_sandbox/modern/saved_figs/10_realistic_sinct.pdf} 
		\caption{Asymptotically approaches 0 at $\pm \infty$}
	\end{figure}
\end{frame}

\begin{frame}{Approach - Mock Microphone Recording}
	Where precisely is the center of chirp in this audio sample?
	\begin{figure}
		\centering
		\includegraphics[width=0.68 \textwidth]{../../python_correlation_sandbox/modern/saved_figs/07_example_envt_sinct.pdf}
%		\includegraphics[width=0.68 \textwidth]{../../python_correlation_sandbox/modern/saved_figs/11_realistic_envt_sinct.pdf}
		\caption{Naive template matching in the time domain is inaccurate}
	\end{figure}
\end{frame}

\begin{frame}[fragile]{Approach - Cross Correlation}
	\begin{figure}
		
		Definition of Cross Correlation for complex signals with period $T$:
		\begin{align}
			(f \star g)(\tau) & = \int_{-\infty}^{+\infty} \overline{f(t)} g(t+\tau) \,dt = \int_{t_0}^{t_0+T} \overline{f(t)} g(t+\tau) \,dt
		\end{align}
		
	\end{figure}
	
	
	\begin{lstlisting}
int cross_correlation(vector sgnlA, vector sgnlB, int offset)
{
	int result = 0;
	for (int i = 0; i < sgnlA.size(); ++i)
	{
		size_t period = sgnlA.size();
		size_t indexB = (i+offset) % period; // pseudocode
		result += std::conj(sgnlA[i]) * sgnlB[indexB];
	}
	return result;
}
	\end{lstlisting}
\end{frame}


\begin{frame}{Cross Correlation}
	Cross-correlating the chirp with the input audio stream produces peaks where signals are "similar"
	\begin{figure}
		\centering
		\includegraphics[width=0.68 \textwidth]{../../python_correlation_sandbox/modern/saved_figs/08_example_correlation_surface_around_expected_peak_500_of_block_period.pdf}
		\caption{Results zoomed and centered around true peak location}
	\end{figure}
\end{frame}

\begin{frame}{Cross Correlation}
	True peak location will always be between 2 audio samples. 
	
	$\pm1$ sample  $\approx \pm 22.6 \mu$s.  range error $\approx \pm 0.3$in. angular error $ \approx \pm6.9^\circ$
	
	\begin{figure}
		\centering
		\includegraphics[width=0.7 \textwidth]{../../python_correlation_sandbox/modern/saved_figs/15_realistic_correlation_surface_around_expected_peak_10_of_block_period.pdf}
		\caption{Results zoomed and centered around true peak location}
	\end{figure}
\end{frame}


\begin{frame}{Solutions that don't work}
	\begin{figure}
		\begin{itemize}
			\item Interpolating the discrete correlation surface's derivative directly is inaccurate
			\item Fitting a 2nd order polynomial to the area around the raw audio is super inaccurate - unclean signal separation
			\item Fitting a 2nd order polynomial to the area around the correlation peak is also inaccurate - but worth investigating
		\end{itemize}
	\end{figure}
\end{frame}

\begin{frame}{Our Solution}
	\begin{figure}
		\begin{itemize}
			\item Re-express the chirp and recorded audio as continuous-time signals using the frequency domain
			\item Use those to express the correlation surface as a continuous-time function instead of a discrete-time one
			\item Find a coarse estimate of the peak using fast, industry standard algorithms (FFT)
			\item Refine the coarse estimate by using a root-solver on the analytic derivative of the continuous-time correlation surface
			\item Profittttt
		\end{itemize}
	\end{figure}
\end{frame}